\pagenumbering{roman}
\setcounter{page}{1}

% Daftar Isi
\tableofcontents
\newpage

% Daftar Tabel
\listoftables
\newpage

% Daftar Gambar
\listoffigures
\newpage

% Tabel Kontribusi
\section*{Tabel Kontribusi}
\addcontentsline{toc}{section}{Tabel Kontribusi}

Berikut disajikan rincian pembagian tanggung jawab dan kontribusi masing-masing anggota kelompok.
\vspace{0.3cm}
\begin{footnotesize}
\begin{spacing}{1.05}
\renewcommand\tabularxcolumn[1]{m{#1}}
\begin{xltabular}{\textwidth}{>{\centering\arraybackslash}m{2.2cm} m{3.8cm} X}
\caption{Tabel Kontribusi pada Milestone 2} \label{tab:kontribusi_m2} \\
\toprule
\textbf{NIM} & \textbf{Nama} & \textbf{Kontribusi} \\
\midrule
\endfirsthead

\multicolumn{3}{l}{\textit{\footnotesize Lanjutan Tabel Kontribusi Milestone 2}} \\
\toprule
\textbf{NIM} & \textbf{Nama} & \textbf{Kontribusi} \\
\midrule
\endhead

\bottomrule
\endfoot

13523061 & Darrel Adinarya Sunanda &
\begin{itemize}[leftmargin=*, noitemsep, topsep=1pt, parsep=0pt]
    \item \textbf{Sharding Storage Multi-Node (P1):} Merancang dan mengimplementasikan arsitektur \textit{database sharding} PostgreSQL 3-node dengan volume \textit{hostPath} node-lokal dan kueri agregasi \textit{scatter-gather}.
    \item \textbf{HA Leader Leasing \& Fencing (P2):} Mengimplementasikan mekanisme \textit{distributed leader election} Aggregator dengan \textit{fencing token} monotonik dan \textit{source checkpoints} untuk mencegah \textit{split-brain} dan \textit{zombie writes}.
    \item \textbf{Audit API \& Verifikasi Integritas:} Mengembangkan \textit{endpoint} audit kanonis (\texttt{/internal/audit/events} dan \texttt{/internal/audit/checksum}) dengan verifikasi SHA-256 dan \textit{keyset pagination}.
    \item \textbf{Downstream Field-Report Consumer (P3):} Mengimplementasikan konsumer topik \texttt{bnpb.field-reports} pada Aggregator dengan deduplikasi \textit{payload hash} dan \textit{transactional outbox}.
    \item \textbf{AI Situation Summarizer (Bonus M2):} Mengembangkan konsumer peringatan bencana berbasis AI \textit{non-blocking} menggunakan LLM lokal (\texttt{qwen3.5:9b}).
    \item \textbf{Unified Observability Stack:} Mengimplementasikan observabilitas terpadu mencakup \textit{dashboard} Grafana, metrik Prometheus antarlayanan, dan \textit{distributed tracing} OpenTelemetry/Tempo dengan propagasi W3C \textit{traceparent}.
    \item \textbf{Infrastruktur \& Koordinasi:} Menyusun \textit{Kubernetes manifests}, \textit{network policies}, dokumen ADR 005--008, serta koordinasi pengujian dan pelaporan.
\end{itemize} \\
\midrule
13524110 & Jennifer Khang &
\begin{itemize}[leftmargin=*, noitemsep, topsep=1pt, parsep=0pt]
    \item \textbf{Gateway Laporan Lapangan (P3):} Mengimplementasikan \texttt{POST /field-reports} pada Client API dengan autentikasi berbasis \textit{role}, respons instan \textit{202 Accepted}, penjaminan idempotensi (\textit{Idempotency-Key}), dan \textit{crash-safe staging} ke Redpanda.
    \item \textbf{Web Analytics Dashboard (P1):} Mengembangkan layanan dan antarmuka dasbor analitik bencana \textit{real-time} dengan agregasi multi-sumber, sinkronisasi delta audit, dan indikator \textit{partial data}.
    \item \textbf{Coastal Sensor Ingestion \& QoS (P5):} Mengembangkan \textit{ingestion poller} sensor pantai (\texttt{sensor-pantai}) dengan \textit{circuit breaker}, \textit{fallback cache}, koordinasi \textit{leaseGuard}, karantina ke DLQ, serta \textit{middleware QoS load shedding}.
    \item \textbf{Unified Observability Stack:} Mengimplementasikan observabilitas terpadu mencakup \textit{dashboard} Grafana, metrik Prometheus antarlayanan, dan \textit{distributed tracing} OpenTelemetry/Tempo dengan propagasi W3C \textit{traceparent}.
    \item \textbf{Pembuatan Laporan:} Berkontribusi dalam penyusunan laporan.
\end{itemize} \\
\midrule
18223113 & Aldoy Fauzan Avanza &
\begin{itemize}[leftmargin=*, noitemsep, topsep=1pt, parsep=0pt]
    \item \textbf{Evolusi Skema \& DLQ Parking (P3):} Mengembangkan simulasi mode perubahan skema inkompatibel pada \texttt{mock-pvmbg}, kebijakan penanganan evolusi skema pada Aggregator, penyimpanan DLQ (\texttt{dlq\_messages}), dan \textit{replay API}.
    \item \textbf{Satellite Wildfire Ingestion (P4):} Merancang dan mengimplementasikan klien \textit{stream} satelit (SSE), algoritma pengelompokan spasio-temporal (\textit{hotspot clustering}), \textit{event-time windowing}, penanganan data terlambat, dan metrik \textit{satellite lag}.
    \item \textbf{Horizontal Pod Autoscaler (Bonus M2):} Mengimplementasikan autoscaling horizontal (HPA) dinamis berbasis utilisasi CPU pada pod Client API dan Aggregator.
    \item \textbf{Pembuatan Laporan:} Berkontribusi dalam penyusunan laporan.
\end{itemize} \\
\end{xltabular}
\end{spacing}
\end{footnotesize}

\newpage
\pagenumbering{arabic}
\setcounter{page}{1}
